
\begin{proposition}[A rational double cover of a smooth tacnode image]\label{cdp:st:C.2.RAT.2}
Retain the torus-fibre, rational-normalization and specified
ruled-boundary morphism of \ref{cdp:st:C.2.RAT.1}. Suppose that $C$ is not
contracted, that $H=h(C)_{\rm red}$ is a smooth rational curve,
that the generic transversal along $H$ is an ordinary tacnode,
and that $C\to H$ has degree two. Then $r\le2$ and
\[
                         e(S)\ge2+N_0.
\tag{C.RA.18}\label{cdp:eq:C.RA.18}
\]
The assertion allows arbitrary vertical conductor and singular
normalization away from the generic points of $C$.
\end{proposition}
\begin{proof}
The two branches of a generic tacnode are both supplied by
$C\to H$; no other strict conductor component can dominate $H$.
Apply the exact normal-jet construction (A.T.5)--(A.T.7) on
$H\simeq\PP^1$. Here the curve conductor divisor $\Xi$ is zero.
The quotient $Q$ of its rank-one Hessian has degree $q_0\le0$.
The quadratic coefficient $a$ and the quartic invariant
$\Delta=ae-b^2/4$ are nonzero sections with
\[
 A:=\deg\operatorname{div}(a)=-2q_0,\qquad
 \deg\operatorname{div}(\Delta)=8+2q_0=8-A.
\tag{C.RA.19}\label{cdp:eq:C.RA.19}
\]
In particular $A$ is a nonnegative even integer. These are
sections of the actual lines $M\otimes Q^{-2}$ and
$(M\otimes Q^{-1}\otimes R^{-2})^2$, with
$\deg M=0$ and $\deg R=-2-q_0$.

The normalized conductor cover is obtained by adjoining
$\sqrt{-\Delta}$, by the quadratic equation for the ratio of
the first normal coordinate to the square of the radical
coordinate. Thus its branch divisor consists exactly of the
odd-order zeros of $\Delta$, as proved in (A.T.7). Since
$C\simeq\PP^1$ and the map has degree two, Riemann--Hurwitz
gives two distinct branch points on $H$. At each of them
$a$ vanishes: a nonzero Hessian would give the constant
transverse $A_3$ product by \ref{cdp:st:A.5.17} and hence an unramified
conductor cover. The coordinate comparison
(A.T.10)--(A.T.12) excludes a simple zero of $\Delta$ at
such a point. Its odd zero order is consequently at least
three. Equations \eqref{cdp:eq:C.RA.19} now give
\[
                        A\ge2,\qquad 8-A\ge6.
\tag{C.RA.20}\label{cdp:eq:C.RA.20}
\]
Therefore $A=2$. The two branch points are exactly the
zeros of $a$, each simple, and the zeros of $\Delta$ have
order three there, with no further zeros.

Every original fibre $F_j$ of (C.RA.1) gives a positive
vertical boundary component meeting $C$. The positivity
is preserved under blowup at its intersection with the
coefficient-two horizontal section, and these $r$
intersection points on $C$ are distinct. At any such point
the Hessian coefficient at its image must vanish. Otherwise
\ref{cdp:st:A.5.17} supplies the constant $A_3$ product, whose two
normalization branches are smooth and have no vertical
conductor; the minimal resolution $\pi$ is an isomorphism
near them. Each of the two zeros of $a$ is a branch point
of $C\to H$, and consequently has just one point above it
on $C$. It follows that $r\le2$.

Finally the exact complement calculation and (A.RT.21) give
\[
 e(S)=2-r+N_0+e(\Gamma_{\rm red})
                  \ge4-r+N_0\ge2+N_0,
\]
which proves \eqref{cdp:eq:C.RA.18}.
\end{proof}


